% !PRINTED 93
\weberpage{93}
%<<CONT
or, since the integers $x_k$ are arbitrary, the sums
\[
\sum_{1,\mu}^{h}\frac{\alpha_h\,\gamma_{k,h}}{c_h}\qquad k=1,2,\ldots \nu
\]
must be integers. If one then sets, according to (15) and (11) \S.\,21,
\[
\gamma_{k,h}=c_{k,h}\,e_{k,h},\quad c_h=\delta_{k,h}\,c_{k,h},\quad i_k=\delta_{k,h}\,i_{k,h},
\]
it follows that
\[
\sum^{h}\frac{\alpha_h\,e_{k,h}}{\delta_{k,h}}=\sum^{h}\frac{\alpha_h\,i_{k,h}\,e_{k,h}}{i_k}
\]
must be integers, and this requirement is equivalent to the system of congruences
\[
\sum^h \alpha_h\,i_{k,h}\,e_{k,h}\equiv 0\ (\mathrm{mod}\ i_k),
\]
or, written more fully,
\weq{3}{\begin{aligned}
&e_{1,1}i_{1,1}\alpha_1+e_{1,2}i_{1,2}\alpha_2+\cdots+e_{1,\mu}i_{1,\mu}\alpha_\mu\equiv 0\ (\mathrm{mod}\ i_1)\\
&e_{2,1}i_{2,1}\alpha_1+e_{2,2}i_{2,2}\alpha_2+\cdots+e_{2,\mu}i_{2,\mu}\alpha_\mu\equiv 0\ (\mathrm{mod}\ i_2)\\
&\cdots\cdots\cdots\cdots\cdots\cdots\cdots\cdots\cdots\cdots\cdots\cdots\cdots\cdots\cdots\cdots\cdots\cdots\cdots\cdots\\
&e_{\nu,1}i_{\nu,1}\alpha_1+e_{\nu,2}i_{\nu,2}\alpha_2+\cdots+e_{\nu,\mu}i_{\nu,\mu}\alpha_\mu\equiv 0\ (\mathrm{mod}\ i_\mu). % [sic: Weber prints i_\mu; the row pattern requires i_\nu.]
\end{aligned}}
By these congruences the numbers $\alpha_1, \alpha_2, \ldots \alpha_\mu$ are, as must be the case, determined only with respect to the moduli $c_1, c_2, \ldots c_\mu$, because $i_{k,h}c_h=i_k c_{k,h}$ is divisible by $i_k$; and \gesperrt{to find the indices of all numbers of the group $\mathfrak{A}$, one has simply to seek out all solutions of the congruences (3).}

\begin{center}
\rule{4cm}{0.4pt}
\end{center}
